% SPDX-License-Identifier: Apache-2.0
\section{A Peripheral Written in Software}
\label{sec:x-remote}

A peripheral is usually hardware before it is anything else, and the
program that drives it is written against a datasheet and debugged
against a simulation. \code{Remote} turns that round. It sits on
AXI-Lite where any peripheral would, and answers nothing itself: each
transaction it accepts leaves on a channel as an \code{Ask}, and an
\code{Answer} comes back on another. What is at the far end of those
two channels is not the peripheral's business, and in this example
it is a wire. \code{RemoteLink} turns an \code{Ask} into one Ethernet
frame and a frame into an \code{Answer}; the Ethernet MAC of
\code{//docs:eth} puts it on the wire and takes it off at the other
end; and a program reads the frame, works out the answer, and
sends a frame back. Nothing in the design knows that the device it
is talking to is twenty lines of Rust.

So a peripheral can be designed as software, against the bus it will
really sit on, with the rest of the system running as it really runs.
When the software is right, hardware replaces it and nothing else in
the design changes.

What the bus sees is a slow device. One transaction remains
outstanding at a time, holding the transaction channel until the
response arrives, which is what AXI-Lite allows and what a program on the
far side of a wire needs. A device that never answers must not be able to
stop the bus, so a transaction unanswered for \code{T} cycles is
answered \code{SlvErr} by the peripheral itself. Each request includes
a tag echoed in the response, ensuring safety: an answer arriving after
the timeout counter expires holds a tag that has already retired, allowing
the peripheral to discard it without misattributing it to subsequent
transactions.

The run does four things, in the order a person would try them. It
writes a word and reads it back, so the program holds the state. It
reads a word the program computes rather than stores, the number of
writes it has seen, which is a line of software and a register file
nobody would build. It reads an address the program has no word for,
and the refusal arrives as \code{SlvErr} on the bus. Then the program
ignores one frame, and the bus is told the device failed after the
thousand cycles the peripheral waits rather than waiting for ever;
then the program answers again and the next read is served.

A transaction is about seven hundred cycles here, because each frame
is stored whole by a MAC before it goes and there are two of them
each way. That is the price of the wire, and it is why the patience
is a parameter: a program in the same simulation answers in two
cycles and one across a network takes thousands.

The same protocol is spoken by \code{//tools/remote}, a program in Go
that serves a device over TCP connections routed through SSH tunnels.
Its test asserts the same frame, byte for byte, that the hardware sends,
so the two implementations are held to one protocol rather than to
each other.

Both units are lowered, and the build simulates both netlists against
this run under \code{nvc} and Verilator.

\begin{figure}[htbp]
\centering
\input{remote_timing.tex}
\caption{A transaction becoming a frame: \code{ask} dispatches the
request from the peripheral, \code{at} counts the twenty-seven frame
bytes driven onto \code{byte}, and \code{tx\_en} asserts when the MAC
transmits the completed frame. \code{busy} remains asserted throughout
transmission and the subsequent round-trip response latency.}
\label{fig:x-remote}
\end{figure}

\lstinputlisting[caption={\code{ex\_remote.rs}. Run with \code{bazel run //lib/examples:ex\_remote}.},label={lst:x-remote}]{lib/examples/ex_remote.rs}

\lstinputlisting[style=ascii,caption={What \code{ex\_remote} printed when the build ran it: the bus and the program, then the lowering.},label={lst:x-remote-out}]{out_remote.txt}
